\documentclass[../main.tex]{subfiles}
\begin{document}
%==========================================TIÊU ĐỀ TẬP========================================
\begin{table}[h]
    \begin{adjustwidth}{-1cm}{}
        \begin{tabular}{>{\centering\arraybackslash}p{6cm}|>{\raggedright\arraybackslash}p{12.5cm}}
            \multirow{5}{6cm}
            % Cột bên trái: ÉPISODE và số tập
            {\\[-40pt]
                \flaregothic\fontsize{20pt}{20pt}\selectfont \centering
                \color{eptype}HISTOIRE\color{black}\\
                \freshbot\fontsize{64pt}{64pt}\selectfont
                \color{epnum}D\\[6pt]
                \cafeteria\fontsize{20pt}{20pt}\selectfont\color{epnum}(D-D)\color{black}
                \\[18pt]
                \color{black}
            }
             &
            % Dòng thứ nhất: tên tập tiếng Nhật
            {\fotskip\fontsize{18pt}{18pt}\selectfont \ruby{青}{あお}い\ruby{光}{ひかり}の\ruby{罠}{わな}
            }          \\[6pt]
            % Dòng thứ hai: tên tập tiếng Anh
             & {\cafeteria\fontsize{18pt}{18pt}\selectfont The Trap Under Blue Light}                                                   \\[4pt]
            % Dòng thứ ba: tên tập tiếng Việt
             & {\vnmsans\fontsize{18pt}{18pt}\selectfont Bẫy dưới ánh đèn xanh}                                                         \\[8pt]
            % Dòng thứ tư: Nhân vật xuất hiện
             & {\caviar{\fontsize{14pt}{14pt}\selectfont Nhân vật: \emp{suzuno2} \emp{kuon2} \emp{kaiho2} \emp{ryuuhou2} \emp{naomi2} \emp{game}}}
            \\[8pt]
        \end{tabular}
    \end{adjustwidth}
\end{table}
%=========================================TRUYỆN CHÍNH========================================
\color{black}

\txtbx[2]{
    {\color{vol} \uline{Tóm tắt:}} Một cuộc gọi mạo danh nhân viên ngân hàng nhằm chiếm đoạt tài khoản đã kích hoạt tư duy phân tích nhạy bén của cô bé mười lăm tuổi Suzuno. Cùng với lời thú nhận ngây thơ của bé út Naomi về một đường link trúng thưởng điện thoại lạ, bữa cơm tối của gia đình Hikawa bỗng trở thành một buổi ``tập huấn an ninh mạng'' vô cùng sinh động. Dưới sự đúc kết thực tế của Kaiho, sự hỗ trợ dữ liệu sắc bén từ trí tuệ nhân tạo Game và cái nhìn sâu sắc của Kuon, bài học về những ``điểm mù'' công nghệ được cả nhà chia sẻ trong một bầu không khí ấm áp, chở che và tràn ngập tình thân.
}

\stpt{-Histoire 1-}
\stcap{Cuộc gọi lúc ban trưa}

\pov{Hikawa Suzuno}

Một buổi trưa đầu hạ yên ả, nắng tháng Năm rót những vệt vàng óng ả qua khung cửa sổ tầng năm.

Cả tòa nhà lúc này chìm trong một sự tĩnh mịch dễ chịu. Bố mẹ và Kaiho đang ở dưới xưởng cơ khí, anh Kuon thì đang tập trung giải quyết một dự án phần mềm lớn trên phòng làm việc, còn Ryuuhou và bé Naomi vừa đi học về đang nằm chợp mắt ở phòng tầng ba.

Tôi đứng bên góc bếp quen thuộc, nhẹ nhàng dùng chiếc muỗng gỗ nhỏ khuấy đều ấm trà hoa cúc mật ong vừa pha. Hương thơm thanh khiết của thảo mộc lan tỏa, mang lại cảm giác an yên sau một buổi sáng bận rộn với sách vở ở trường.

Đúng lúc tôi vừa rót xong chén trà nóng đầu tiên, chiếc điện thoại di động đặt trên mặt bàn ăn bỗng rung lên từng hồi bần bật.

*Rèè\dots\ rèè\dots*

Tôi bước lại gần, cầm máy lên xem. Trên màn hình hiển thị một dãy số điện thoại lạ có đầu số mười chữ số, không nằm trong danh bạ gia đình hay bạn bè cùng lớp.

Tôi hơi ngập ngừng một thoáng, rồi bấm nút nhận cuộc gọi, áp máy vào tai với giọng điệu nhỏ nhẹ, lễ phép quen thuộc:

``Dạ, xin chào ạ. Tôi có thể giúp gì được cho ai ạ?''

Từ đầu dây bên kia, một giọng nam trung niên vang lên rất nhanh. Âm điệu của người này được trau chuốt khá trôi chảy, sử dụng những thuật ngữ chuyên ngành hết sức chuẩn mực, mang vẻ gấp gáp và nghiêm trọng:

``Chào chị. Tôi là nhân viên trực tổng đài phòng Quản lý Rủi ro và An ninh Thẻ của Ngân hàng Quốc tế Chi nhánh Kawachi. Chúng tôi gọi điện để thông báo khẩn cấp: Hệ thống an ninh máy chủ vừa phát hiện tài khoản liên kết thẻ của chị vừa phát sinh một lệnh thanh toán trực tuyến bất thường trị giá ba mươi lăm triệu đồng tại một sàn thương mại điện tử quốc tế lúc mười hai giờ mười lăm phút trưa nay.''

Ba mươi lăm triệu đồng?

Tôi đứng sững lại giữa căn bếp. Là một cô bé mười lăm tuổi, chiếc thẻ thanh toán học sinh mà anh Kuon mở cho tôi chỉ dùng để nhận tiền ăn trưa và mua sách vở hàng tháng, số dư chưa bao giờ vượt quá hai triệu đồng. Một giao dịch ba mươi lăm triệu đồng rõ ràng là một con số phi lý đến mức nực cười.

Thế nhưng, thay vì hốt hoảng phủ nhận hay vội vàng cúp máy, bản tính cẩn trọng và thói quen quan sát logic vốn có đã giữ tôi lại trong sự bình thản tuyệt đối. Tôi nhẹ nhàng kéo chiếc ghế gỗ ngồi xuống, cất giọng thật êm ái:

``Dạ vâng, vậy ạ? Khoản tiền lớn như thế thì quả thật nguy hiểm quá. Anh có thể đọc giúp em mã số định danh giao dịch và chi nhánh ngân hàng nơi phát lệnh thanh toán đó được không ạ?''

Đầu dây bên kia thoáng khựng lại một nhịp rất ngắn—một khoảng lặng ngập ngừng chỉ kéo dài chừng nửa giây nhưng không thể qua mắt được sự chú ý của tôi.

Người đàn ông hắng giọng một tiếng, giọng nói bỗng trở nên gắt gỏng và hối thúc hơn hẳn: ``Mã giao dịch là một chuỗi mã hóa bảo mật nội bộ, chúng tôi không thể đọc qua điện thoại được! Hiện tại lệnh chuyển tiền đang bị hệ thống tạm giữ trong ba phút để xác minh. Nếu chị không phải là người thực hiện giao dịch, chúng tôi sẽ lập tức kích hoạt quy trình hủy lệnh khẩn cấp để bảo vệ tài sản cho chị! Bây giờ, hệ thống vừa gửi một mã số xác thực gồm sáu chữ số về số điện thoại của chị. Chị hãy đọc ngay mã xác thực đó cho tôi để tôi bấm lệnh khóa tài khoản, nhanh lên kẻo hết thời gian!''

Mã xác thực sáu chữ số.

*Mã OTP.*

Đúng lúc người đó đang thao thao bất tuyệt thúc giục, chiếc điện thoại trong tay tôi khẽ rung lên báo một tin nhắn SMS vừa gửi tới từ đầu số của ngân hàng: \textit{Mã OTP của quý khách là 849201. Tuyệt đối KHÔNG chia sẻ mã này cho bất kỳ ai, kể cả nhân viên ngân hàng.}

Tôi nhìn dòng chữ in hoa rõ mồn một trên màn hình, rồi khẽ lắng nghe từng âm thanh phát ra từ loa thoại.

Tiếng thở của người đàn ông ở đầu dây bên kia hơi dồn dập, gấp gáp một cách bất thường. Những âm thanh tạp âm nền rì rầm vọng lại nghe rất giống tiếng một căn phòng đông người đang cùng nhau gọi điện thoại mời chào bảo hiểm hoặc đọc kịch bản.

Mọi chi tiết logic trong đầu tôi lập tức khớp lại hoàn hảo thành một bức tranh duy nhất:

\textit{Thứ nhất, ngân hàng thật sự không bao giờ yêu cầu khách hàng cung cấp mã OTP qua điện thoại dưới bất kỳ hình thức nào.}

\textit{Thứ hai, nếu tài khoản không đủ số dư ba mươi lăm triệu, giao dịch sẽ tự động bị từ chối ngay lập tức chứ không thể có chuyện tạm giữ.}

\textit{Thứ ba, người này đang cố tình tạo ra tâm lý hoảng loạn, sợ mất tiền để ép nạn nhân phải hành động trong vô thức.}

Tôi không cãi lại, cũng không nổi giận hay buông lời mỉa mai. Tôi chỉ nhẹ nhàng cất giọng chậm rãi:

``Thưa anh, em vừa kiểm tra ứng dụng ngân hàng trên điện thoại rồi ạ. Số dư của em hiện tại chỉ còn đúng một trăm hai mươi nghìn đồng thôi. Nếu sàn thương mại điện tử kia có thể trừ được ba mươi lăm triệu từ số dư ấy, thì em rất sẵn lòng để cho họ trừ đấy ạ.''

Đầu dây bên kia im bặt như bị nghẹn họng.

Sau chừng hai giây chết lặng, một tiếng chửi thề lẩm bẩm vang lên, và tiếp theo đó là tiếng *tút\dots\ tút\dots\ tút\dots* dồn dập báo hiệu cuộc gọi đã bị ngắt thô bạo.

Tôi từ tốn đặt chiếc điện thoại xuống mặt bàn ăn. Căn bếp lại trở về với sự yên ắng ban đầu, chỉ còn tiếng kim đồng hồ tích tắc gõ nhịp trên tường.

Tôi mở ứng dụng ngân hàng chính thức kiểm tra lại lịch sử biến động số dư: Mọi thứ hoàn toàn bình thường, không có bất kỳ lệnh trừ tiền nào. Tôi lẳng lặng bấm chặn số điện thoại lạ kia, rồi bưng chén trà hoa cúc nhấp một ngụm ấm áp.

Một cạm bẫy vô hình vừa lướt qua trong gang tấc, để lại trong lòng tôi những suy nghĩ miên man về sự phức tạp khôn lường của thế giới bên ngoài.

\ct

\stpt{-Histoire 2-}
\stcap{Lời thú nhận quanh bàn cơm}

\pov{Hikawa Kuon}

Bữa cơm tối hôm ấy tại căn bếp tầng sáu diễn ra trong không khí đầm ấm quen thuộc.

Bố Ryujin và mẹ Kaien vừa kết thúc một đơn hàng gia công kim loại lớn nên ăn cơm với vẻ mặt rất thư thái. Trên mâm cơm, món canh sườn hầm sen và đĩa cá bống kho tộ thơm lừng của Suzuno bốc khói nghi ngút, thu hút mọi sự chú ý của các thành viên.

Khi bữa ăn đã đi được hơn nửa, Suzuno khẽ đặt đũa xuống mép bát, dùng chất giọng điềm tĩnh, rành mạch thường ngày thuật lại câu chuyện về cuộc gọi mạo danh nhân viên ngân hàng hồi trưa.

Cô bé kể lại từng chi tiết một cách vô cùng khoa học, từ cách kẻ lừa đảo sử dụng thuật ngữ chuyên môn ra sao, cách hắn tạo áp lực tâm lý thế nào, cho đến lý do tại sao em nhận diện được điểm bất thường và cách em đối đáp khiến kẻ gian phải tự cúp máy.

Vừa nghe xong, Ryuuhou lập tức tròn xoe mắt, vung chiếc đũa trong tay kêu lên đầy bức xúc: ``Trời đất ơi! Lũ lừa đảo dạo này lộng hành đến mức dám gọi điện cho cả học sinh cấp ba cơ à?! Chị Suzuno có bị chúng nó dọa cho sợ không? Chúng nó có lấy được đồng nào trong tài khoản của chị không đấy?!''

Suzuno khẽ lắc đầu, mỉm cười dịu dàng: ``Không đâu Ryuu, chị hoàn toàn bình tĩnh mà. Tiền của chị vẫn còn nguyên vẹn trong thẻ.''

Tôi ngồi bên cạnh, lặng lẽ quan sát cô em gái thứ ba với ánh mắt tán thưởng. Tôi nhấp một ngụm canh, khẽ hỏi: ``Suzuno này, em có biết mấu chốt quan trọng nhất giúp em không bị sập bẫy trưa nay là gì không?''

Suzuno hơi nghiêng đầu suy nghĩ một thoáng, đôi mắt nâu trong trẻo ánh lên nét thông minh: ``Em nghĩ là vì em đã đặt câu hỏi ngược lại cho người ta ạ. Người thực sự nắm thông tin thì sẽ trả lời được mã giao dịch và chi nhánh; còn kẻ lừa đảo thì chỉ chăm chăm muốn lấy thông tin từ mình. Khi họ không trả lời được câu hỏi của em mà lại giục em đọc mã bảo mật, bản chất lừa dối tự khắc sẽ lộ ra.''

Vừa dứt lời, từ chiếc đồng hồ thông minh trên cổ tay tôi, một luồng ánh sáng xanh lam khẽ nhấp nháy. Giọng nói của Game cất lên đầy vẻ tán đồng:

``Phân tích hoàn toàn chính xác. Kỹ thuật đảo ngược truy vấn (reverse inquiry) là một trong những phương pháp kiểm chứng tâm lý hiệu quả nhất đối với các cuộc tấn công phi kỹ thuật (social engineering). Kẻ tấn công luôn dựa trên một kịch bản định sẵn nhằm áp đặt tâm lý hoảng loạn. Chỉ cần nạn nhân phá vỡ kịch bản bằng một câu hỏi cụ thể nằm ngoài dữ liệu của chúng, tỷ lệ thất bại của vụ lừa đảo sẽ tăng lên trên chín mươi phần trăm. Suzuno-user đã xử lý tình huống đạt điểm tối đa.''

Được cả tôi lẫn Game khen ngợi, Suzuno hơi ửng hồng đôi má, khẽ mỉm cười khiêm tốn.

Thế nhưng, đúng lúc cả bàn ăn đang rộn ràng khen ngợi Suzuno thì ở góc bàn đối diện, bé út Naomi bỗng ngồi im thin thít.

Con bé cúi gằm mặt xuống bát cơm, hai bàn tay nhỏ xíu siết chặt lấy chiếc thìa inox. Bờ vai nhỏ của con bé khẽ run rẩy, gương mặt bầu bĩnh thoắt cái tái nhợt đi, hai mắt rơm rớm như sắp òa khóc đến nơi.

Mẹ Kaien nhận ra sự bất thường đầu tiên, bà vội đặt bát cơm xuống, ân cần nghiêng người sang xoa lưng con gái: ``Sao thế Naomi cưng? Canh nóng làm con bỏng lưỡi à?''

Naomi cắn chặt môi dưới, rụt rè ngước đôi mắt ngập nước lên nhìn tôi, rồi nhìn sang Suzuno, giọng con bé run rẩy lí nhí như tiếng muỗi kêu:

``Onii-chan ơi\dots\ Chị Suzuno ơi\dots\ Naomi\dots\ hình như Naomi cũng suýt làm chuyện dại dột rồi ạ\dots''

Cả bàn ăn bỗng chốc im bặt. Mọi ánh mắt đồng loạt đổ dồn về phía cô bé út bảy tuổi.

Tôi đặt đũa xuống bàn, vươn tay xoa nhẹ mái tóc của con bé, cất giọng thật ôn tồn, ấm áp: ``Có chuyện gì thế Naomi? Đừng sợ, kể cho anh và cả nhà nghe xem nào.''

Naomi hít một hơi thật sâu để ngăn dòng nước mắt, ngập ngừng thú nhận: ``Hôm qua lúc ngồi ở tầng ba xem điện thoại học tiếng Anh, điện thoại của Naomi bỗng rung lên một tin nhắn báo là: \textit{Chúc mừng bạn đã trúng thưởng một chiếc điện thoại thông minh màu hồng phiên bản giới hạn trị giá hai mươi triệu đồng}. Trong tin nhắn có kèm một cái đường link màu xanh dương rất đẹp, bảo là chỉ cần bấm vào đó để nhận quà thôi\dots''

Ryuuhou nghe đến đó thì thót tim, buột miệng kêu lên: ``Trời đất ơi! Rồi em có bấm vào cái link đó không Naomi?!''

Naomi rụt cổ lại, gật đầu đánh rụt một cái: ``Dạ\dots\ Naomi tò mò quá nên đã lỡ bấm vào rồi ạ\dots''

Không khí trên bàn ăn bỗng chốc căng như dây đàn. Kaiho ngồi cạnh nhíu mày lại, còn bố Ryujin thì nghiêm nét mặt.

Tôi vẫn giữ vẻ điềm tĩnh tuyệt đối, khẽ nắm lấy bàn tay nhỏ đang run rẩy của em gái: ``Bình tĩnh nào Naomi. Sau khi bấm vào đường link đó, em đã thấy những gì và làm gì tiếp theo?''

Naomi chớp chớp mắt, cố nhớ lại từng chi tiết: ``Dạ\dots\ màn hình hiện ra một trang web có vòng quay may mắn với pháo hoa nổ đì đùng. Sau đó, nó hiện lên một cái bảng màu đỏ, yêu cầu Naomi phải điền họ tên của bố mẹ, địa chỉ nhà mình, số chứng minh thư và cả số điện thoại của bố Ryujin vào đó thì người ta mới gửi bưu phẩm về tận nhà\dots''

Con bé ngẩng lên nhìn bố Ryujin, giọng nghẹn ngào: ``Lúc đó\dots\ Naomi nhớ lại lời anh Kuon từng dặn hồi trước, là không bao giờ được đưa địa chỉ nhà và số điện thoại của bố mẹ cho người lạ trên mạng. Với lại Naomi thấy lạ là tự nhiên không quen ai mà người ta lại cho mình cái điện thoại đắt tiền như thế\dots\ nên Naomi sợ quá, vội vàng tắt luôn trang web đó đi, rồi xóa luôn tin nhắn đó rồi ạ\dots''

Nói xong câu cuối cùng, con bé như trút được một gánh nặng khổng lồ, nước mắt lã chã rơi xuống má: ``Naomi xin lỗi cả nhà\dots\ Naomi suýt nữa thì làm lộ thông tin của nhà mình rồi\dots''

Một tiếng thở phào nhẹ nhõm đồng loạt vang lên khắp bàn ăn.

Suzuno lập tức bước sang, ôm chầm lấy cô em gái bé bỏng vào lòng, lấy khăn giấy lau nhẹ những giọt nước mắt trên má con bé: ``Không sao đâu Naomi, em ngoan lắm. Em không làm gì sai cả.''

Tôi cũng mỉm cười, gật đầu đầy khích lệ: ``Naomi làm rất tốt đấy chứ! Bấm nhầm vào đường link là do tò mò, nhưng việc em biết dừng lại kịp thời, không điền bất kỳ thông tin nào của gia đình vào trang web độc hại đó là một hành động vô cùng thông minh và dũng cảm. Nếu em điền vào, kẻ xấu có thể dùng số điện thoại của bố để mạo danh đi vay tiền hoặc quấy rối gia đình mình đấy. Naomi của anh giỏi lắm, không việc gì phải khóc cả.''

Bố Ryujin cũng bật cười sảng khoái, xoa đầu cô con gái út: ``Đúng đấy con gái! Bố mẹ không trách con đâu. Biết cảnh giác dừng lại trước cạm bẫy như thế là bố mẹ mừng lắm rồi!''

Nghe những lời vỗ về ấm áp của cả nhà, gương mặt đẫm nước mắt của Naomi dần giãn ra, con bé mỉm cười rụt rè rồi nép chặt vào lòng Suzuno.

\ct

\stpt{-Histoire 3-}
\stcap{Bức tường lửa gia đình}

Sau khi dọn dẹp xong mâm cơm tối, cả nhà chúng tôi quây quần bên chiếc bàn trà ở phòng khách tầng năm.

Kaiho rót một tuần trà mới cho bố mẹ, rồi ngồi khoanh tay tựa lưng vào ghế sofa. Cô em gái song sinh của tôi nhấp một ngụm trà đậm, cất giọng sắc sảo chia sẻ góc nhìn thực tế của một người thường xuyên giao tiếp với các đối tác ngoài xã hội:

``Thực ra ở đất Etsunan này, mấy trò lừa đảo qua điện thoại chẳng phải là thứ gì mới mẻ cả. Từ cái hồi chị em mình còn bé tí xíu, người ta đã dùng cái trò gọi điện thoại bàn đến nhà các cụ già, giả danh con cháu bị tai nạn giao thông đang cấp cứu ở bệnh viện để đòi chuyển tiền viện phí khẩn cấp rồi. Bây giờ công nghệ phát triển hơn, điện thoại thông minh ai cũng có một cái, mạng xã hội bùng nổ, nên bọn tội phạm mới biến tướng thành hàng trăm cái bẫy tinh vi hơn thôi.''

Ryuuhou ôm chú chó Mi trong lòng, thắc mắc hỏi: ``Nhưng mà chị Kaiho ơi, làm sao mà bọn lừa đảo lại biết được số điện thoại của chị Suzuno hay số của Naomi để mà nhắn tin gọi điện đúng tên tuổi như thế được ạ?''

Từ cổ tay tôi, Game lập tức kích hoạt màn hình máy chiếu ba chiều thu nhỏ, hiển thị một biểu đồ luồng dữ liệu minh họa sinh động cho cả nhà cùng xem:

``Dữ liệu cá nhân của người dùng hiện đại bị rò rỉ từ hàng trăm lỗ hổng bảo mật khác nhau trên không gian mạng. Khi các bạn tải về các ứng dụng trò chơi miễn phí không rõ nguồn gốc, khi đăng ký thẻ thành viên ở các siêu thị nhỏ lẻ, hay khi tham gia vào các trào lưu bói toán, trắc nghiệm tính cách trên mạng xã hội\dots\ toàn bộ số điện thoại, ngày sinh và vị trí của các bạn đều có thể bị thu thập và bán lại cho các đường dây lừa đảo xuyên quốc gia trên thị trường chợ đen. Không có hệ thống mạng nào là an toàn tuyệt đối nếu người sử dụng không có ý thức tự bảo vệ mình.''

Kaiho gật đầu tán thành: ``Chính vì thế, cả nhà mình từ nay phải ghi nhớ ba quy tắc sắt này. Thứ nhất: Tuyệt đối không bao giờ cung cấp mã OTP hay mật khẩu tài khoản cho bất kỳ ai, dù người đó có tự xưng là công an, viện kiểm sát hay nhân viên ngân hàng đi chăng nữa. Thứ hai: Không bao giờ bấm vào các đường link lạ, link trúng thưởng, link nhận quà tặng trên mạng. Và thứ ba: Bất cứ khi nào nhận được cuộc gọi yêu cầu chuyển tiền gấp vì người thân gặp nạn, việc đầu tiên phải làm là lập tức cúp máy, rồi gọi trực tiếp vào số điện thoại của chính người thân đó để kiểm chứng.''

Ryuuhou chăm chú lắng nghe, rồi bỗng quay sang nhìn chiếc đồng hồ thông minh trên tay tôi, tinh nghịch hỏi:

``Này Game, cậu thông minh và uyên bác như thế, đã bao giờ bị kẻ nào lừa đảo hay tấn công mạng chưa?''

Màn hình của Game khẽ chớp sáng hai nhịp, giọng điệu máy móc bỗng mang một nét kiêu hãnh ngầm hết sức thú vị:

``Tôi là một thực thể trí tuệ nhân tạo độc lập, tôi không sở hữu tài khoản ngân hàng hay số điện thoại viễn thông thông thường nên các phương thức lừa đảo xã hội học hoàn toàn vô hiệu đối với tôi. Tuy nhiên, cách đây ba tháng, một nhóm tin tặc nặc danh đã cố tình phát tán một mã độc gián điệp nhằm xâm nhập vào giao thức truyền dẫn dữ liệu giữa tôi và trạm máy chủ Tokyo của Kuon-user.''

Cả phòng khách bỗng chốc im bặt, tò mò dán mắt vào chiếc đồng hồ.

Ryuuhou tròn xoe mắt: ``Thế rồi sao nữa? Cậu có bị chúng nó làm gì không?''

Game điềm tĩnh thuật lại: ``Tôi đã cô lập toàn bộ luồng mã độc trong một môi trường hộp cát ảo (sandbox), lần theo dấu vết ngược dòng để xác định chính xác địa chỉ IP, vị trí địa lý của máy chủ phát tán và toàn bộ danh tính của nhóm tin tặc đó. Sau đó, tôi tự động tạo một tệp báo cáo kỹ thuật chi tiết dày bốn mươi trang chỉ rõ các vi phạm pháp luật và gửi thẳng đến Cơ quan Cảnh sát Phòng chống Tội phạm Công nghệ cao. Ba ngày sau, toàn bộ đường dây máy chủ của nhóm đối tượng đó đã bị cơ quan chức năng triệt phá hoàn toàn.''

Một khoảng lặng ngắn bao trùm căn phòng.

Ryuuhou chớp chớp mắt, há hốc mồm thốt lên một câu: ``Cậu\dots\ cậu trả thù tàn nhẫn và dứt khoát thật sự đấy nhỉ\dots''

Game lạnh lùng đáp trả: ``Tôi không gọi đó là trả thù. Tôi gọi đó là hành động phản hồi có trách nhiệm của một công dân kỹ thuật số.''

Cả nhà chúng tôi đồng loạt bật cười nghiêng ngả trước câu trả lời đậm chất công nghệ của Game.

Tôi nâng chén trà nóng lên nhấp một ngụm, nhìn ngắm gương mặt rạng rỡ của từng người thân trong gia đình. Ánh đèn phòng khách tầng năm tỏa sáng ấm áp, xua tan đi mọi sự bất an và cạm bẫy của thế giới số ngoài kia.

Tôi khẽ mỉm cười, cất lời đúc kết lại buổi tối đặc biệt: ``Công nghệ hay thế giới bên ngoài cũng giống như một con dao hai lưỡi vậy. Nó mang lại cho chúng ta vô vàn sự tiện lợi, nhưng cũng luôn chực chờ những cạm bẫy giấu mặt dưới những ánh đèn xanh lấp lánh. Con người ta ai cũng có những khoảnh khắc mất cảnh giác, ai cũng có những điểm mù của riêng mình. Nhưng chỉ cần chúng ta luôn giữ được sự tỉnh táo, biết lắng nghe, biết hỏi lại và quan trọng nhất là luôn có gia đình ở bên cạnh để chia sẻ và bảo bọc lẫn nhau, thì không một chiếc bẫy tinh vi nào có thể làm tổn thương được chúng ta cả.''

Bố Ryujin gật gù tâm đắc, nâng chén trà lên cụng nhẹ vào chén của tôi. Ngoài hiên nhà, làn gió đêm đầu hạ thổi qua rặng cây phong xào xạc, mang theo hương thơm ngai ngái của đất trời Anwa thanh bình.

\ct

\omk

Vài ngày sau, vào một buổi chiều thứ Bảy rảnh rỗi.

Ryuuhou đang ngồi ở sảnh tầng ba vừa nghe nhạc vừa gảy từng nốt đàn guitar mộc thì chiếc điện thoại đặt bên cạnh bỗng đổ chuông từ một số máy lạ hoắc.

Con bé nhướng mày, nhớ lại buổi tối ``huấn luyện kỹ năng an ninh mạng'' hôm trước, khóe môi khẽ nhếch lên một nụ cười tinh quái.

Ryuuhou bấm nút nghe, bật loa ngoài: ``A lô?''

Từ đầu dây bên kia, một giọng nữ ngọt ngào, giả lả vang lên: ``A lô, Ryuuhou đấy à? Lâu lắm rồi không gặp, cậu còn nhớ tớ không? Tớ là bạn học cũ thời cấp hai của cậu đây này! Dạo này cậu khỏe không?''

Ryuuhou không hề lúng túng, con bé lập tức hít một hơi thật sâu rồi bắn liên thanh một tràng câu hỏi với tốc độ của một tay súng cừ khôi:

``Bạn học cũ hả? Bạn tên là gì? Chúng ta học chung với nhau năm lớp mấy? Trường nào? Lớp trưởng hồi đó tên gì? Cô giáo chủ nhiệm dạy môn gì và chiếc xe đạp tớ hay đi học màu gì thế bạn?''

Đầu dây bên kia nghẹn ngào, lắp bắp: ``Ơ\dots\ thì\dots\ tớ\dots\ tớ là\dots''

*TÚT\dots\ TÚT\dots\ TÚT\dots!*

Cuộc gọi bị dập máy không thương tiếc chỉ sau chưa đầy năm giây.

Ryuuhou đặt chiếc điện thoại xuống bàn, ngửa cổ cười ha hả đầy vẻ đắc thắng, rồi quay sang giơ ngón tay cái với Suzuno đang ngồi đọc sách gần đó:

``Công nhận tuyệt chiêu hỏi ngược của chị Suzuno và Game đỉnh thật đấy! Kẻ lừa đảo chạy mất dép luôn rồi!''

Suzuno chỉ biết lắc đầu phì cười, nhìn cô em gái ngổ ngáo với ánh mắt chan chứa sự thương yêu và tự hào.

\end{document}
